\documentclass[a4paper,12pt]{article}

% -----------------------------------------
% Pakete und Einstellungen (Projektstil)
% -----------------------------------------
\usepackage[a4paper, left=2.5cm, right=2.5cm, top=2.5cm, bottom=2.5cm]{geometry}
\usepackage[ngerman]{babel}
\usepackage{fontspec}
\setmainfont{Arial}
\usepackage{csquotes}
\usepackage{graphicx}
\usepackage{float}
\usepackage{hyperref}
\usepackage{listings}
\usepackage{xcolor}
\usepackage{color}
\usepackage{setspace}
\usepackage{booktabs}
\usepackage{multirow}
\usepackage[table]{xcolor}
\usepackage{array}
\usepackage{adjustbox}
\usepackage{pifont}   % fuer \ding{51} (Haken)

% -----------------------------------------
% Farben und Listing-Stil
% -----------------------------------------
\definecolor{lightgray}{rgb}{0.95,0.95,0.95}
\definecolor{darkgray}{rgb}{0.3,0.3,0.3}
\definecolor{codegreen}{rgb}{0.0,0.5,0.0}
\definecolor{codepurple}{rgb}{0.5,0.0,0.5}
\definecolor{tableheadcolor}{rgb}{0.2,0.4,0.6}

\lstset{
    backgroundcolor=\color{lightgray},
    basicstyle=\ttfamily\small,
    breaklines=true,
    breakatwhitespace=false,
    frame=single,
    rulecolor=\color{darkgray},
    keywordstyle=\color{codepurple}\bfseries,
    stringstyle=\color{codegreen},
    commentstyle=\color{darkgray}\itshape,
    showstringspaces=false,
    tabsize=4,
    captionpos=b
}

% Python-Stil
\lstdefinestyle{python}{
    language=Python,
    morekeywords={None,True,False,str,int,ValueError},
}

% Log-Stil (kein Syntax-Highlighting, aber einheitlich)
\lstdefinestyle{log}{
    basicstyle=\ttfamily\footnotesize,
    backgroundcolor=\color{lightgray},
    frame=single,
    breaklines=true,
    breakatwhitespace=false,
}

% Tabellenkopf-Makro
\newcommand{\tableheader}[1]{%
    \rowcolor{tableheadcolor}\color{white}\textbf{#1}%
}

\setstretch{1.5}

% -----------------------------------------
% Dokument
% -----------------------------------------
\begin{document}

% =========================================
% Titelbereich
% =========================================
\begin{center}
    {\Large \textbf{Testprotokoll -- TaskGrid Distributed System}}\\[0.5em]
    {\large Verteilte Systeme}\\[1em]
    \begin{tabular}{ll}
        \textbf{Datum:}               & 08. Juni 2026 \\
        \textbf{System:}              & TaskGrid (Verteilte Systeme) \\
        \textbf{Testumgebung:}        & Docker Compose (\texttt{docker compose up -d --build}) \\
        \textbf{Protokoll-Referenz:}  & Strukturierte Logs (Kap.~13) -- Dispatcher-Log und Worker-Client-Log \\
    \end{tabular}
\end{center}

\vspace{1em}
\hrule
\vspace{1.5em}

% =========================================
\section*{Fehlertests}
% =========================================

% =========================================
\newpage
\section*{F-01 -- Worker-Container wird während der Bearbeitung gekillt}
% =========================================

\begin{table}[H]
    \centering
    \begin{adjustbox}{max width=\textwidth}
    \begin{tabular}{>{\bfseries}p{4.5cm} p{10cm}}
        \toprule
        Feld & Inhalt \\
        \midrule
        Test-Nr              & F-01 \\
        Fehlerfall           & Worker-Ausfall während der Bearbeitung (\texttt{docker kill worker-wait}) \\
        Ausgangssituation    & System läuft mit Dispatcher, Namensdienst und \texttt{worker-wait}.
                               Dispatcher-Timeout = 3\,s, Max-Retries = 2.
                               Ein \texttt{wait}-Task mit 30 Sekunden Laufzeit wird gestartet. \\
        Durchgeführte Aktion & Task wird an den Worker dispatcht.
                               Während der Bearbeitung wird der Worker-Container mittels
                               \texttt{docker kill} abrupt beendet. \\
        Erwartetes Verhalten & Der Task darf nicht erfolgreich abgeschlossen werden.
                               Dispatcher erkennt den Ausfall über Timeout oder gRPC-Fehler und setzt den
                               Task auf \texttt{FAILED}. Der Namensdienst entfernt den Worker nach
                               ausbleibenden Heartbeats. \\
        Tatsächliches Verhalten & Der Task endet im Status \texttt{FAILED}.
                               Der Dispatcher bleibt weiterhin funktionsfähig.
                               Andere Tasks können erfolgreich verarbeitet werden.
                               Der Namensdienst markiert den Worker als offline. \\
        Bewertung            & \textbf{PASS} -- Der Worker-Ausfall wird korrekt erkannt.
                               Der Task wird nicht fälschlich abgeschlossen und das System bleibt stabil. \\
        \bottomrule
    \end{tabular}
    \end{adjustbox}
    \caption{Fehlertest F-01: Worker-Ausfall während der Bearbeitung}
\end{table}

\subsection*{Relevante Logs}

\begin{lstlisting}[style=log]
DISPATCH_TIMEOUT_fired
DISPATCH_worker_unreachable
DISPATCH_stale_rpc_error_ignored
\end{lstlisting}

\subsection*{Screenshot}

Terminal-Ausgabe von \texttt{docker kill worker-wait} sowie Dispatcher- und Namensdienst-Logs.

% =========================================
\newpage
\section*{F-02 -- Unbekannter Tasktyp gesendet}
% =========================================

\begin{table}[H]
    \centering
    \begin{adjustbox}{max width=\textwidth}
    \begin{tabular}{>{\bfseries}p{4.5cm} p{10cm}}
        \toprule
        Feld & Inhalt \\
        \midrule
        Test-Nr              & F-02 \\
        Fehlerfall           & Unbekannter Tasktyp (\texttt{unknown\_type}) \\
        Ausgangssituation    & System läuft mit \texttt{worker-sum}, \texttt{worker-reverse}
                               und \texttt{worker-hash}.
                               Für den Typ \texttt{unknown\_type} existiert kein registrierter Worker. \\
        Durchgeführte Aktion & Client sendet einen \texttt{POST\_TASK}-Request mit
                               \texttt{task\_type="unknown\_type"} und Payload \texttt{"test"}. \\
        Erwartetes Verhalten & Dispatcher akzeptiert den Task zunächst und vergibt eine Task-ID.
                               Beim Dispatch findet der Namensdienst keinen passenden Worker.
                               Der Task darf niemals den Zustand \texttt{COMPLETED} erreichen. \\
        Tatsächliches Verhalten & Der Namensdienst liefert keine Worker zurück.
                               Der Dispatcher protokolliert den Fehler und belässt den Task in
                               \texttt{QUEUED}, \texttt{RETRYING} oder \texttt{FAILED}.
                               Das System bleibt stabil und verarbeitet weiterhin andere Tasks. \\
        Bewertung            & \textbf{PASS} -- Das System behandelt unbekannte Tasktypen korrekt,
                               ohne abzustürzen oder fehlerhafte Ergebnisse zu erzeugen. \\
        \bottomrule
    \end{tabular}
    \end{adjustbox}
    \caption{Fehlertest F-02: Unbekannter Tasktyp}
\end{table}

\subsection*{Relevante Logs}

\begin{lstlisting}[style=log]
POST_TASK_accepted
LOOKUP_WORKER_no_workers
DISPATCH_no_worker
\end{lstlisting}

\subsection*{Screenshot}

Dispatcher-Log mit Meldung über fehlende Worker für den Tasktyp \texttt{unknown\_type}.

% =========================================
\newpage
\section*{F-03 -- Worker antwortet nicht innerhalb des Timeouts}
% =========================================

\begin{table}[H]
    \centering
    \begin{adjustbox}{max width=\textwidth}
    \begin{tabular}{>{\bfseries}p{4.5cm} p{10cm}}
        \toprule
        Feld & Inhalt \\
        \midrule
        Test-Nr              & F-03 \\
        Fehlerfall           & Worker-Timeout während der Bearbeitung \\
        Ausgangssituation    & Dispatcher-Timeout = 3\,s, Max-Retries = 2.
                               Ein \texttt{wait}-Worker verarbeitet einen Task mit
                               einer Laufzeit von 30 Sekunden und antwortet daher nicht rechtzeitig. \\
        Durchgeführte Aktion & Ein Task mit \texttt{wait=30} wird an den Worker gesendet.
                               Der Dispatcher überwacht die Bearbeitung mittels Timeout-Timer. \\
        Erwartetes Verhalten & Nach Ablauf des Dispatcher-Timeouts wird der Task erneut eingeplant.
                               Nach Erreichen der maximalen Anzahl von Wiederholungsversuchen
                               (\texttt{MAX\_RETRIES}) wird der Task endgültig auf \texttt{FAILED} gesetzt. \\
        Tatsächliches Verhalten & Der Task durchläuft die Zustände
                               \texttt{DISPATCHED \textrightarrow{} TIMEOUT \textrightarrow{} RETRYING \textrightarrow{} FAILED}.
                               Anschließend können neue Tasks weiterhin erfolgreich verarbeitet werden. \\
        Bewertung            & \textbf{PASS} -- Der Timeout-Mechanismus funktioniert korrekt.
                               Der Task wird nach den definierten Wiederholungsversuchen beendet und das
                               System bleibt weiterhin stabil. \\
        \bottomrule
    \end{tabular}
    \end{adjustbox}
    \caption{Fehlertest F-03: Worker-Timeout}
\end{table}

\subsection*{Relevante Logs}

\begin{lstlisting}[style=log]
DISPATCH_TIMEOUT_fired
TIMEOUT_retrying
TIMEOUT_max_retries_exceeded
\end{lstlisting}

\subsection*{Screenshot}

Dispatcher-Log mit Timeout-Ereignissen und Statuswechseln des Tasks.

% =========================================
\newpage
\section*{Gesamtzusammenfassung}
% =========================================

\begin{table}[H]
    \centering
    \begin{adjustbox}{max width=\textwidth}
    \begin{tabular}{lll}
        \toprule
        \textbf{Test} & \textbf{Fehlerfall} & \textbf{Ergebnis} \\
        \midrule
        F-01 & Worker-Absturz während der Bearbeitung & \ding{51}~PASS \\
        F-02 & Unbekannter Tasktyp                    & \ding{51}~PASS \\
        F-03 & Worker-Timeout                         & \ding{51}~PASS \\
        \bottomrule
    \end{tabular}
    \end{adjustbox}
    \caption{Ergebnisse der Fehlertests}
\end{table}

\subsection*{Beobachtungen}

\textbf{Worker-Ausfall korrekt erkannt (F-01):} Der Dispatcher erkennt den Ausfall über Timeout
oder gRPC-Fehler und setzt den Task auf \texttt{FAILED}. Der Namensdienst markiert den Worker
als offline nach ausbleibenden Heartbeats. Das System bleibt in allen Fällen stabil.

\textbf{Unbekannte Tasktypen sicher behandelt (F-02):} Der Namensdienst liefert beim Lookup
keine Worker zurück (\texttt{LOOKUP\_WORKER\_no\_workers}). Der Dispatcher protokolliert den
Fehler korrekt, ohne abzustürzen oder fehlerhafte Ergebnisse zu erzeugen.

\textbf{Timeout-Mechanismus funktional (F-03):} Nach Ablauf des konfigurierten Timeouts wird
der Task erneut eingeplant. Nach Erreichen von \texttt{MAX\_RETRIES} wird der Task endgültig
auf \texttt{FAILED} gesetzt. Neue Tasks können anschließend weiterhin erfolgreich verarbeitet werden.

Alle drei Fehlerfälle wurden erfolgreich getestet. In keinem der getesteten Szenarien kam es
zu einem Systemabsturz oder zu inkonsistenten Taskzuständen.

\end{document}